% ═══════════════════════════════════════════════════════════════
% Section 6: Neural Impact Chain
% ═══════════════════════════════════════════════════════════════

\section{Neural Impact Chain}
\label{sec:nic}

The Neural Impact Chain (NIC) is a six-stage threat-modeling method that associates
attack techniques with DSM-5-TR and ICD-10-CM diagnostic category
references~\cite{apa2022dsm5tr}. The codes are \textbf{diagnostic category references for
threat modeling, not diagnostic claims}: an association says that a technique and a
condition category share an address and a function, not that the technique causes the
condition. The question the NIC helps a threat modeler ask is \emph{which clinical
categories should a review of this technique consider?} The mappings have not been reviewed
by clinicians (Section~\ref{sec:limitations}).

\subsection{Pipeline Architecture}

The NIC traces each technique through six stages:

\begin{enumerate}
  \item \textbf{Technique}: The TARA attack technique (e.g., cortical signal
        injection)
  \item \textbf{Address}: The code(s) the technique is tagged at (e.g., N7
        Neocortex, N6 Allocortex; version 1 memberships in the snapshot)
  \item \textbf{Neural Structure}: The anatomical structure at risk (e.g.,
        hippocampus, prefrontal cortex, amygdala)
  \item \textbf{Cognitive Function}: The function that structure supports (e.g.,
        memory consolidation, executive control, emotional regulation)
  \item \textbf{NISS Score}: The neural impact score, particularly the BI, CR,
        CD, and NP metrics that correlate with clinical outcomes
  \item \textbf{Category reference}: The DSM-5-TR/ICD-10-CM code(s) for the
        condition category most closely associated with disruption of that function
\end{enumerate}


\subsection{NISS-to-DSM Bridge}

A rule assigns each technique to a DSM-5-TR category cluster according to which NISS
metric dominates its profile. The rule is a modeling choice, not a clinical finding:

\begin{itemize}
  \item \textbf{BI-driven} $\to$ Motor/Neurocognitive cluster: movement and
        neurocognitive categories
  \item \textbf{CR/CD-driven} $\to$ Cognitive/Psychotic cluster: perceptual and
        cognitive categories
  \item \textbf{CV-driven} $\to$ Mood/Trauma cluster: trauma- and stressor-related
        categories
  \item \textbf{NP/RV-driven} $\to$ Persistent/Personality cluster: categories of
        lasting change
  \item \textbf{No neural impact} $\to$ Non-diagnostic: Silicon-only
        techniques with no category association
\end{itemize}

\subsection{Coverage Statistics}

All \SnapshotTechniques{} TARA techniques have been mapped through the NIC pipeline:

\begin{table}[H]
\centering
\caption{Neural Impact Chain mapping results across \SnapshotTechniques{} techniques.}
\label{tab:nic-stats}
\small
\begin{tabular}{@{}l r@{}}
\toprule
\textbf{Metric} & \textbf{Value} \\
\midrule
Techniques mapped                 & 103 / \SnapshotTechniques{} (94.5\%) \\
Unique category codes             & 29 \\
Diagnostic clusters               & 5 \\
Direct association                 & 52 (47.7\%) \\
Indirect association               & 9 (8.3\%) \\
No association (silicon-only)      & 42 (38.5\%) \\
\bottomrule
\end{tabular}
\end{table}

\subsection{Diagnostic Cluster Distribution}

The five diagnostic clusters and their technique counts:

\begin{table}[H]
\centering
\caption{DSM-5-TR diagnostic category clusters (snapshot), for threat modeling.}
\label{tab:clusters}
\small
\begin{tabularx}{\textwidth}{@{}l r X@{}}
\toprule
\textbf{Cluster} & \textbf{Count} & \textbf{Representative DSM-5-TR Codes} \\
\midrule
Non-diagnostic          & 42 & --- (silicon-only techniques) \\
Mood/Trauma             & 21 & F43.10 (PTSD), F32.9 (MDD), F44.9 \\
Cognitive/Psychotic     & 17 & F06.0 (psychosis), R41.3 (cognitive decline) \\
Motor/Neurocognitive    & 16 & G25.9 (movement disorder), G31.84 \\
Persistent/Personality  &  7 & F07.0 (personality change) \\
\bottomrule
\end{tabularx}
\end{table}

The Mood/Trauma cluster is the largest diagnostic category cluster (21 techniques),
because the bridge rule sends consent-dominated techniques there. That is a consequence of
the rule, not an observation about outcomes in people.

\subsection{Association Classes}

Each technique is placed in one of three association classes:

\begin{itemize}
  \item \textbf{Direct} (52 techniques, 47.7\%): the technique acts on the
        structure or function associated with the category. Example: unauthorised
        cortical stimulation is associated with seizure-related categories.
  \item \textbf{Indirect} (9 techniques, 8.3\%): the association runs through a
        downstream effect. Example: exfiltration of neural data is associated with
        anxiety-related categories through the effect of the disclosure, not through
        the neural signal.
  \item \textbf{None} (42 techniques, 38.5\%): Silicon-only techniques with no
        direct neural interaction and thus no category association.
\end{itemize}

\subsection{Example Walkthrough}

Consider \textbf{QIF-T0001, Signal injection}, as stored in the live registry. (The v1.5
walkthrough gave a vector, a score of 8.7 and three codes that did not match the registry
entry; it is replaced here.)

\begin{enumerate}
  \item \textbf{Technique}: QIF-T0001 (signal injection at the electrode-tissue interface)
  \item \textbf{Address}: I0, N1
  \item \textbf{Mechanism}: electrical current delivered at the interface, modulating
        local field potentials
  \item \textbf{NISS}: \texttt{BI:H/CR:H/CD:H/CV:E/RV:P/NP:T} $\to$ 6.1 (Medium), PINS
        set because BI is High
  \item \textbf{Category references}: F43.2, F45, F44.4 (primary); F82 (secondary)
  \item \textbf{Association class}: direct
\end{enumerate}

\noindent Its evidence status is theoretical: the technique carried \emph{confirmed} until
2026-10-16, when its sources were found not to demonstrate signal injection
(Section~\ref{sec:tara}).

\subsection{Impact Chains as Associations (Live Dataset)}

The live impact-chains dataset (5,546 rows) is built by joins: each technique to every
structure sharing one of its addresses, then to every pathway touching that structure, then
to that pathway's condition codes. A row therefore records a shared address, not that the
technique reaches the structure. The join currently yields rows for 91 of the
\LiveTechniques{} techniques and 10 structures. A technique with no row has no association
recorded, which is not evidence that it has no effect. Because the joins run on addresses,
the move to addressing version 2 (Section~\ref{sec:hourglass}) changes them; the snapshot
figures above use version 1.

